% !TEX root = task-4-3-4.tex
% Editable task bank. Append new BankTask blocks at the end of the document.
\documentclass[12pt,a4paper]{article}
\usepackage{physics-textbook}
\usepackage{graphicx,needspace}
\usetikzlibrary{decorations.pathmorphing}
\newif\ifShowSolutions
\ifdefined\TutorEdition
  \ShowSolutionstrue
\else
  \ShowSolutionsfalse % Change to \ShowSolutionstrue to display solutions.
\fi
\newcounter{banktask}
\newsavebox{\BankTaskBox}
\newcommand{\BankTaskContent}[3]{%
  \noindent\textbf{\thebanktask.}\quad #2\par
  \ifShowSolutions
    \Needspace{50mm}% Keep the solution heading with its opening text or diagram.
    #3\par
    \ifdefined\TutorEdition
      \ifcsname TutorSource#1\endcsname
        \nopagebreak\smallskip{\small\csname TutorSource#1\endcsname\par}%
      \fi
    \fi
  \fi
}
\newcommand{\BankTask}[3]{%
  \refstepcounter{banktask}\par\medskip
  \sbox{\BankTaskBox}{\begin{minipage}{\linewidth}%
    \BankTaskContent{#1}{#2}{#3}\end{minipage}}%
  \ifdim\dimexpr\ht\BankTaskBox+\dp\BankTaskBox\relax>\dimexpr\textheight-20mm\relax
    \Needspace{35mm}\BankTaskContent{#1}{#2}{#3}%
  \else
    \Needspace{\dimexpr\ht\BankTaskBox+\dp\BankTaskBox\relax}%
    \noindent\usebox{\BankTaskBox}\par
  \fi
  \medskip
}
\begin{document}
\BookSubsection{4.3.4}{Полное внутреннее отражение}


% Source sdamgia-1924: https://phys-ege.sdamgia.ru/problem?id=1924
\BankTask{sdamgia-1924}{
Синус предельного угла полного внутреннего отражения на границе стекло  — воздух равен $\frac{8}{13} .$ Чему равен абсолютный показатель преломления стекла? (Ответ округлить до сотых.)
}{
\smallskip\noindent\textbf{Решение}\par
Абсолютный показатель преломления связан с синусом угла полного внутреннего отражения соотношением:

\[\frac{1}{n} = \frac{\sin \alpha _{\,\text{полн}}. \,\text{вн}. \,\text{отр}.}{1} .\]

Следовательно:

\[n= \frac{1}{\sin \alpha _{\,\text{полн}}. \,\text{вн}. \,\text{отр}.} = \frac{13}{8} \approx 1{,}63.\]

Ответ: $1{,}63$.
}

% Source sdamgia-3488: https://phys-ege.sdamgia.ru/problem?id=3488
\BankTask{sdamgia-3488}{
Чему равен синус предельного угла полного внутреннего отражения при переходе света из вещества с $n=1{,}5$ в вещество с $n=1{,}2$?
}{
\smallskip\noindent\textbf{Решение}\par
\begin{center}\begin{tikzpicture}[book diagram,>=stealth]\draw(-2.3,0)--(2.3,0);\draw(0,-1.8)--(0,1.2);\draw[thick,->](-2.00000,-1.5)--(0,0);\draw[thick,->](0,0)--(2.2,0);\node at(1.4,-.9){$n_1$};\node at(1.4,.8){$n_2$};\draw(0,-.55)arc(-90:-143.13010:.55);\node at(-.45,-1.25){$\alpha_{\text{пред}}$};\end{tikzpicture}\end{center}

Явление полного внутреннего отражения наблюдается при переходе света из среды оптически более плотной в среду оптически менее плотную. Предельный угол падения луча определяется тем условием, что преломленный луч идет параллельно границе раздела двух сред, то есть синус угла преломления равен $1$. Поэтому, согласно закону Снеллиуса, для синуса предельного угла полного внутреннего отражения имеем:

\[\frac{\sin \alpha}{1} = \frac{n_2}{n_1} = \frac{1{,}2}{1{,}5} =0{,}8.\]

Ответ: $0{,}8$.
}

% Source sdamgia-3496: https://phys-ege.sdamgia.ru/problem?id=3496
\BankTask{sdamgia-3496}{
Чему равен синус угла полного внутреннего отражения при переходе света из вещества, где скорость света равна $0{,}5c,$ в вещество, где скорость света равна $0{,}8 c$? (c  — скорость света в вакууме.)
}{
\smallskip\noindent\textbf{Решение}\par
\begin{center}\begin{tikzpicture}[book diagram,>=stealth]\draw(-2.3,0)--(2.3,0);\draw(0,-1.8)--(0,1.2);\draw[thick,->](-1.20096,-1.5)--(0,0);\draw[thick,->](0,0)--(2.2,0);\node at(1.4,-.9){$n_1$};\node at(1.4,.8){$n_2$};\draw(0,-.55)arc(-90:-128.68219:.55);\node at(-.45,-1.25){$\alpha_{\text{пред}}$};\end{tikzpicture}\end{center}

Скорость света в среде связана с ее абсолютным показателем преломления и скоростью света в вакууме $\,\text{с}$ соотношением $n v =c.$ Следовательно, чем меньше скорость света в веществе, тем больше ее показатель преломления, тем более оптически плотной она является. Явление полного внутреннего отражения наблюдается при переходе света из среды оптически более плотной в среду оптически менее плотную. Предельный угол падения луча определяется тем условием, что преломленный луч идет параллельно границе раздела двух сред, то есть синус угла преломления равен $1$. Поэтому, согласно закону Снеллиуса, для синуса предельного угла полного внутреннего отражения имеем:

\[\frac{\sin \alpha}{1} = \frac{n_2}{n_1} = \frac{v _1}{v _2} = \frac{0{,}5c}{0{,}8c} =0{,}625.\]

Ответ: $0{,}625$.
}

% Source sdamgia-3502: https://phys-ege.sdamgia.ru/problem?id=3502
\BankTask{sdamgia-3502}{
Свет идет из вещества с показателем преломления n в вакуум. Предельный угол полного внутреннего отражения равен $30 ^\circ.$ Чему равен n?
}{
\smallskip\noindent\textbf{Решение}\par
Коэффициент преломления показывает, во сколько раз скорость света в среде меньше, чем скорость света в вакууме, поэтому коэффициент преломления вакуума равен $1$. Явление полного внутреннего отражения наблюдается при переходе света из среды оптически более плотной в среду оптически менее плотную. Предельный угол падения луча определяется тем условием, что преломленный луч идет параллельно границе раздела двух сред, то есть синус угла преломления равен $1$. Поэтому, согласно закону Снеллиуса, для искомого коэффициента преломления имеем:

\[\frac{\sin \alpha _{\,\text{пред}}}{1} = \frac{1}{n} \Longleftrightarrow n= \frac{1}{\sin \alpha _{\,\text{пред}}} = \frac{1}{\sin 30 ^\circ} =2.\]

Ответ: $2$.
}

% Source sdamgia-31414: https://phys-ege.sdamgia.ru/problem?id=31414
\BankTask{sdamgia-31414}{
Синус предельного угла полного внутреннего отражения на границе вода  — воздух равен $0{,}75$. Чему равна скорость света в воде (в $\text{км/с}$)?
}{
\smallskip\noindent\textbf{Решение}\par
По закону преломления света при переходе из воды в воздух $\frac{\sin \alpha}{\sin \gamma} = \frac{v}{c} .$ При полном внутреннем отражении $\sin \gamma =1.$ Тогда: 

\[v =c \sin \alpha _0=300 000 \cdot 0{,}75=225 000 \,\text{км/с}.\]

Ответ: $225$ $000$.
}

% Source sdamgia-8719: https://phys-ege.sdamgia.ru/problem?id=8719
\BankTask{sdamgia-8719}{
Точечный источник света находится в емкости с жидкостью и опускается вертикально вниз от поверхности жидкости. При этом на поверхности жидкости возникает пятно, в пределах которого лучи света от источника выходят из жидкости в воздух. Глубина погружения источника (расстояние от поверхности жидкости до источника света), измеренная через равные промежутки времени, а также соответствующий радиус светлого пятна представлены в таблице. Чему равен показатель преломления жидкости? (с точностью до сотых.)

\begin{center}\begin{tabular}{|c|c|c|c|c|c|c|c|}\hline
Глубина погружения, см&$10$&$20$&$30$&$40$&$50$&$60$&$70$\\\hline
Радиус пятна, см&$12$&$24$&$36$&$48$&$60$&$72$&$84$\\\hline
\end{tabular}\end{center}
}{
\smallskip\noindent\textbf{Решение}\par
Тангенс предельного угла полного внутреннего отражения равен отношению радиуса пятна к глубине погружения. Для всех экспериментальных данных $\tan \alpha _{\,\text{пред}}= \frac{12}{10} =1{,}2.$
Показатель преломления жидкости равен $n= \frac{1}{\sin \alpha _{\,\text{пред}}} = \sqrt{1 + \frac{1}{\tan ^2 \alpha _{\,\text{пред}}}} = \sqrt{1+\frac{1}{1{,}2^2}} \approx1{,}30.$

Ответ: $1{,}30$.
}

% Source sdamgia-3156: https://phys-ege.sdamgia.ru/problem?id=3156
\BankTask{sdamgia-3156}{
Луч света падает на границу раздела «стекло  — воздух». Как изменятся при увеличении показателя преломления стекла следующие три величины: длина волны света в стекле, угол преломления, угол полного внутреннего отражения?

$1$.  Увеличится.
$2$.  Уменьшится.
$3$.  Не изменится.

\begin{center}\begin{tabular}{|c|c|c|}\hline
Длина волны 
света в стекле&Угол 
преломления&Угол полного 
внутреннего отражения\\\hline
&&\\\hline
\end{tabular}\end{center}
}{
\smallskip\noindent\textbf{Решение}\par
При переходе светового пучка из стекла в воздух частота электромагнитных колебаний в световой волне не изменяется, поскольку она не зависит от того, в какой среде распространяется волна. Длина волны связана с частотой электромагнитных колебаний и скоростью распространения соотношением $\lambda \nu =c.$ При увеличении показателя преломления скорость распространения световой волны в среде уменьшается, а значит, уменьшается и длина волны света в стекле.
Согласно закону преломления Снеллиуса, синусы углов падения и преломления при выходе света из стекла в воздух связаны с показателем преломления стекла соотношением $\frac{\sin \alpha _{\,\text{пад}}}{\sin \gamma _{\,\text{прел}}} = \frac{1}{n} .$ Следовательно, при увеличении показателя преломления угол преломления будет увеличиваться.
Наконец, угол полного внутреннего отражения определяется соотношением $\sin \alpha _{\,\text{полн}}.\,\text{внутр}.\,\text{отраж}= \frac{1}{n} .$ Таким образом, увеличение n приведет к уменьшению угла полного внутреннего отражения.

Ответ: $212$.
}

% Source sdamgia-3022: https://phys-ege.sdamgia.ru/problem?id=3022
\BankTask{sdamgia-3022}{
Бассейн глубиной $3\,\text{м}$ заполнен водой, относительный показатель преломления на границе воздух  — вода $1{,}33$. Каков радиус светового круга на поверхности воды от электрической лампы на дне бассейна?
}{
\smallskip\noindent\textbf{Решение}\par
\begin{center}\begin{tikzpicture}[book diagram,>=stealth]\draw(-1,0)--(4.1,0);\fill(0,-2)circle(1.5pt)node[below]{$A$};\fill(0,0)circle(1.5pt)node[above left]{$B$};\fill(2.28084,0)circle(1.5pt)node[above left]{$C$};\draw[thick,->](0,-2)--(0,.85);\draw[thick,->](0,-2)--(.5,0);\draw[thick,->](.5,0)--(.84,1);\draw[dashed](.5,-.8)--(.5,1.1);\draw[thick,->](0,-2)--(2.28084,0);\draw[thick,->](2.28084,0)--(3.9,0);\draw[thick,->](2.28084,0)--(3.42126,-1);\draw[dashed](2.28084,-.9)--(2.28084,.9);\draw(0,-1.45)arc(90:41.2463:.55);\node[right]at(.25,-1.4){$\alpha$};\draw(2.28084,.35)--(2.63084,.35)--(2.63084,0);\node[above right]at(2.38084,.5){$\beta$};\end{tikzpicture}\end{center}

Рассмотрим ход лучей, исходящих из точки А на дне бассейна. Вертикальный луч АВ не изменяет своего направления после прохождения границы раздела, все остальные лучи испытывают преломление или отражаются от границы раздела.
Полное внутреннее отражение происходит начиная с такого значения угла падения $\alpha ,$ при котором угол преломления $\beta$ равен $90 ^\circ.$ Следовательно, предельное значение угла $\alpha ,$ при котором свет выходит из воды, определяется условием $\sin \alpha = \frac{1}{n} .$
Из треугольника АВС найдем радиус светового круга ВС:

\[BC=AB \cdot \tan \alpha = \frac{AB}{\sqrt{n ^2 - 1}} = \frac{3 \,\text{м}}{\sqrt{1{,}33 ^2 - 1}} \approx 3{,}4\,\text{м}.\]

Ответ: $3{,}4\,\text{м}$.
}

% Source sdamgia-3049: https://phys-ege.sdamgia.ru/problem?id=3049
\BankTask{sdamgia-3049}{
У самой поверхности воды в реке летит комар, стая рыб находится на расстоянии $2\,\text{м}$ от поверхности воды. Каково максимальное расстояние до комара, на котором он еще виден рыбам на этой глубине? Относительный показатель преломления света на границе воздух  — вода равен $1{,}33$.
}{
\smallskip\noindent\textbf{Решение}\par
\begin{center}\begin{tikzpicture}[book diagram,>=stealth]\draw(-1.8,0)--(4.5,0);\draw[<->](-.7,-1.5)--(-.7,0)node[midway,left]{$H$};\draw(1.71063,-1.5)--(1.71063,1.3);\fill(1.71063,0)circle(2pt)node[above right]{$K$};\fill(0,-1.5)circle(2pt)node[below right]{$A$};\draw[thick,->](2.627,1.2)--(1.1,0);\draw[thick,->](1.1,0)--(0,-1.5);\draw[thick,->](1.71063,0)--(0,-1.5);\draw[thick,->](4.3,-1.5)--(2.15,0);\draw[thick,->](2.15,0)--(0,-1.5);\draw[dashed](0,-1.5)--(0,0);\draw(0,-.95)arc(90:41.2463:.55);\node at(.3,-.5){$\alpha_{\text{пред}}$};\draw[fill=black!8](-.2,-1.6)..controls(-.8,-1.3)and(-1.3,-1.4)..(-1.45,-1.6)..controls(-1,-1.9)and(-.5,-1.9)..cycle;\draw(-1.45,-1.6)--(-1.7,-1.3)--(-1.7,-1.9)--cycle;\end{tikzpicture}\end{center}

Рыба видит комара, если существует световой луч от него, который, переломившись на границе раздела воздух  — вода, попадет ей в глаз. Вода  — среда оптически более плотная чем воздух, поэтому угол преломления всегда меньше, чем угол падения. Так как комар находится над самой водой, максимальное расстояние между комаром и рыбой определяется предельным углом преломления, который можно найти при помощи закона преломления Снеллиуса: $\frac{1}{\sin \alpha _{\,\text{пред}}} = \frac{n}{1} .$ Из рисунка видно, что расстояние между комаром и рыбой равно

\[AK= \frac{H}{\cos \alpha _{\,\text{пред}}} = \frac{H}{\sqrt{1 - \sin ^2 \alpha _{\,\text{пред}}}} = \frac{H}{\sqrt{1 - 1/n ^2}} = \frac{Hn}{\sqrt{n ^2 - 1}} = \frac{2\,\text{м} \cdot 1{,}33}{\sqrt{1{,}33^2-1}} \approx 3\,\text{м}.\]

Ответ: $3\,\text{м}$.
}

% Source sdamgia-52052: https://phys-ege.sdamgia.ru/problem?id=52052
\BankTask{sdamgia-52052}{
Точечный источник света находится на дне бассейна с плоским дном и прозрачной жидкостью с показателем преломления равным n  =  $1{,}5$. Найдите максимальное время, затраченное светом на прохождение слоя жидкости с последующим выходом в воздух, если минимальное время равно $t_{min}$  =  $50\,\text{нс}$. Сделайте рисунок, на котором постройте ход двух лучей, соответствующих случаям минимального и максимального времени выхода лучей в воздух.
}{
\smallskip\noindent\textbf{Решение}\par
\begin{center}\begin{tikzpicture}[book diagram,>=stealth]\fill[black!7](-1,0)rectangle(3.8,1.7);\draw(-1,0)--(3.8,0)(-1,1.7)--(3.8,1.7);\foreach\x in{-1,-.8,...,3.7}\draw(\x,0)--++(.15,-.15);\fill(0,0)circle(2pt)node[below]{$S$};\draw[thick,->](0,0)--(0,2.3);\node[left]at(0,.8){1};\draw[thick,->](0,0)--(1.52053,1.7);\draw[thick,->](1.52053,1.7)--(3.3,1.7);\draw[dashed](1.52053,.9)--(1.52053,2.3);\draw[<->](-.65,0)--(-.65,1.7)node[midway,left]{$H$};\draw(0,.55)arc(90:48.1897:.55);\node at(.35,.8){$\alpha_0$};\draw(1.52053,1.15)arc(-90:-131.8103:.55);\node at(1.17053,.9){$\alpha_0$};\node at(2.5,.8){$n$};\node at(.95,.55){2};\end{tikzpicture}\end{center}

Точечный источник света находится в более плотной среде, чем воздух. Лучи света, падающего на границу раздела сред под углами большими, чем предельный угол полного внутреннего отражения, не смогут выйти из жидкости. Поэтому луч $1$, затративший на прохождение слоя жидкости с последующим выходом в воздух минимальное время, проходит границу сред без преломления. Луч $2$, затративший на прохождение слоя жидкости с последующим выходом в воздух максимальное время, падает на границу под предельным углом полного внутреннего отражения.
Предельный угол $\alpha_0$ полного внутреннего отражения находим из закона преломления света, соответствующего углу преломления $90^\circ$: $\sin \alpha _0 = \frac{1}{n} .$
Минимальное время на прохождение слоя жидкости затратит луч $1$, идущий перпендикулярно границе раздела сред, а максимальное  — луч $2$, который распространяется под углом $\alpha_0$. Тогда 

$t_{\min} = \frac{H}{V} ,$
$t_{\max} = \frac{H}{V \cos \alpha _0} ,$

где H  — глубина бассейна, V  — скорость света в жидкости.
Имеем: 

\[t_{\max} = \frac{t_{\min} n}{\sqrt{n ^2 - 1}} = \frac{50 \cdot 10 ^{ - 9 } \cdot 1{,}5}{\sqrt{1{,}5 ^2 - 1}} \approx 67 \,\text{нс}.\]

Ответ: $t_{\max} \approx 67 \,\text{нс}.$
}

\end{document}
